%% notes-tex/database/database-expansion-hosts/sections/fill.tex
%% [[experiments.database.expansion-hosts]]
%% https://github.com/Mjvolk3/torchcell/tree/main/notes-tex/database/database-expansion-hosts/sections/fill.tex
%% SOURCE: experiments/database/scripts/build_hosts_candidate_datasets_table.py

\section{The fill pass\secstatus{todo}}\label{sec:fill}

A ranking of 140 rows against a target of 200 was closed on 2026-11-24 by searching five
named veins rather than by widening the existing ones. The pool now holds 201 candidates. Six
sweeps ran: 3-hydroxypropionic acid in \org{P. kudriavzevii} and in \org{S. cerevisiae}, the
LASER strain-engineering database for both of its organisms, and the unserved high-throughput
literature of \org{S. cerevisiae} and \org{E. coli}. Each was given the 170 served datasets and
the 443 rows the three earlier passes had already ranked, so every row below is new to all of
them.

\subsection{3-hydroxypropionic acid, and the structure of its literature\secstatus{todo}}\label{sec:fill-3hp}

Thirteen rows came from one product line across two hosts, and the finding that matters is not
any of them. Of roughly twenty \org{S. cerevisiae} 3-HP papers whose full text was read, four
release machine-readable per-strain numbers and the rest release bar heights. Two independent
sweeps of that literature reached the same count.

The first row to build is Qin et al. 2024, 71 strains running \texttt{QLW1} through
\texttt{QLW71} with no gaps, whose supplementary data carries per-replicate values and a
cumulative genotype string for every strain. Two earlier passes missed it because its title
says yeast rather than \org{Saccharomyces}. One parser trap travels with the row: strain
\texttt{QLW2} is written \texttt{QLW2-HCO3-} in the file, so a strain name matched exactly
rather than as a prefix drops records silently.

Three rows record things the ranking could not absorb. Tong et al. 2021 is a THIRD 3-HP route,
through oxaloacetate rather than beta-alanine or malonyl-CoA, so it needs its own schema leaf;
it is indexed in neither PubMed nor Europe PMC, because its journal is in neither, which is how
a whole pathway stayed invisible to two passes. Tan et al. 2026 in \org{P. kudriavzevii} is
absent from both indexes for the same reason and surfaced only through a DOI registry, with its
abstract null in two more, so nothing about its size is established. And three patents carry
real per-strain tables, the largest giving 11 named strains with both titer and yield, which
the served vocabulary cannot cite at all (Sec.~\ref{sec:pipeline}, Table~\ref{tab:schema}).

Two corrections to the earlier ranking came out of the same pass. The Illinois Data Bank
deposit this wave had ranked as the host's highest-yield 3-HP record is not a phenotype dataset:
it holds a technoeconomic model and genome-scale model code, and the measurements belong to
Tan et al. 2026, which is now ranked on its own identity. And this host has two production
lineages rather than one, SD108 and NRRL~Y-27441, which no row had distinguished.

\subsection{LASER is an index, not a store\secstatus{todo}}\label{sec:fill-laser}

The named source for the fill contributed no rows of its own and three rows by pointing. LASER
was counted from the bytes of its only live release: 456 curated papers, 638 designs, 441
\org{E. coli} and 197 \org{S. cerevisiae}, every file dated 2016-04-27, and the web interface
its paper names no longer resolves. Titer is populated on 181 of 197 yeast designs and yield on
58, there is no rate field at all, and no replicates, no sample count, no uncertainty, no
sequence-level edit and no strain accession. Its own total field sums to 3,472,162 strains built
against 638 curated, because it keeps champion strains and discards the populations they came
from. Every number in it is re-extracted from a primary paper, so it is excluded as a derivative
and used the way its design invites, as a worklist.

What it pointed at is worth more than it is. Kim et al. 2013 measured the anaerobic fermentation
products of 472 single-gene knockouts across five compounds, roughly 2,360 values, behind which
LASER curated seven strains; the table holding them is unlocated, because the publisher refuses
automated fetch and no open copy exists. Valle et al. 2015 screened 150 knockouts in triplicate
for hydrogen, ethanol and glycerol and tabulated all of it in three supplementary files, from
which LASER kept the 12 confirmed hits.

\subsection{What the two unserved-literature sweeps found\secstatus{todo}}\label{sec:fill-unserved}

Forty-four rows came from sweeping the \org{S. cerevisiae} and \org{E. coli} literature for
releases the database does not hold, with every count in the largest rows obtained by
downloading the file and counting its rows rather than from an abstract.

Four releases are larger than the papers that announce them. One \org{E. coli} promoter library
ships 233,469 promoters against a headline of 112,955; another's deposited table carries four
times the measurements its abstract reports; a metabolome time course releases roughly 240 times
the figure its abstract gives, because the abstract counts injections and the release counts
values; and one expression library's abstract reports only its converged subset rather than all
29,249 constructs. Pinning the release rather than the paper is what makes these rows countable,
and it is the reason the ranking quantity is measured and not quoted.

The aggregators were checked and both failed, for opposite reasons. MaveDB holds 60 yeast-target
score sets out of 3,005, and the largest block of them turns out to be in-vitro folding
measurements on constructs optimized for a different organism's codon usage, so they are not
perturbations of a yeast cell. ALEdb has grown to 3,721,969 mutation instances from 13,425
isolates, roughly 94 times what its paper reported, and releases no phenotype number of any
kind, so it is a genotype source. The large yeast variant-effect data lives in transposon
saturation screens, barcoded lineage libraries and non-promoter reporter libraries deposited one
paper at a time.

One shape recurs often enough to be a schema question. Eleven of the twenty largest yeast rows
carry no gene identifier anywhere in their bulk data, because the record key is a lineage
barcode, a raw measured sequence, or a residue substitution. A barcoded lineage is the hardest
of those: the perturbation is unknown by construction and is identified only by its tag, so the
honest representation is a typed absence on the genotype and not a mapping to a gene the release
never names (Table~\ref{tab:schema}).

\subsection{What the fill cost the ranking\secstatus{todo}}\label{sec:fill-cost}

Applying this document's own rule to the new rows moves every pseudomonad fitness row out of the
recommended fifty. The \org{E. coli} libraries are regulatory DNA, so they enter the band that
is applied before tier, and nine of them rank above the Fitness Browser rows on measurements.
The recommended fifty is now 25 \org{S. cerevisiae}, 9 \org{E. coli}, 8 \org{P. kudriavzevii}
and 8 \org{P. putida}, and \org{P. fluorescens}, \org{P. syringae}, \org{P. simiae},
\org{S. stutzeri} and \org{P. aeruginosa} hold none of it.

That outcome contradicts two things this document says. The row count is described as coming
from the pseudomonads, and the per-host floor exists so that every host in the wave starts
(Sec.~\ref{sec:floor}). It also leaves the regulatory-DNA band holding 34 of the 50, at which
point the band no longer separates anything.

The ordering is left as the rule produces it, and the three ways to resolve it are stated rather
than chosen, because which one is right depends on what the wave is for and that is not a
ranking question. The floor could extend to the four thin pseudomonad hosts at one row each,
which keeps the reason the floor was written. The band could take a cap, which is a change to
the selection rule and larger than a fill warrants. Or the ordering stands and four hosts wait
for a later pass.
